% !TEX root = ../Document.tex
%% Abstract
%%
\chapter*{ABSTRACT}\thispagestyle{headings}
\addcontentsline{toc}{compteur}{ABSTRACT}
%
\begin{otherlanguage}{english}

The integration of generative policies such as Flow Matching offers an effective framework for reactive and multimodal trajectory generation in robotic manipulation using visual inputs. However, learned generative models do not provide formal safety guarantees and may produce control commands that lead to collisions in environments containing obstacles not seen during training. Conversely, reactive filters based on Control Barrier Functions provide forward invariance under explicit regularity and execution assumptions, but traditionally rely on simplified geometric representations of the robot and its surroundings.

This thesis presents a decoupled two-layer architecture for the safe real-time execution of trajectories generated by a frozen Flow Matching policy without policy retraining. The framework separates task planning from safety enforcement. The perception system combines the SAM~3 semantic segmentation model with a persistent point-cloud map that preserves obstacle representations during occlusions caused by wrist-mounted camera motions. The nominal policy generates joint-space velocity references at 5~Hz. A reactive safety shield evaluates signed distance fields learned with Bernstein polynomials for the protected robot bodies. For a fixed TopK selection and pointwise distances of class $\mathcal{C}^2$, the soft-min barrier is of class $\mathcal{C}^2$ and is used consistently in value and gradient. The resulting Control Barrier Function constraints are solved at 50~Hz by quadratic programming. Under the stated regularity and execution assumptions, the continuous-time formulation certifies forward invariance of the learned soft-min safe set. Exact-mesh distances are evaluated \textit{a posteriori} on every recorded cycle to verify physical clearance. A reflex brake running at 1~kHz dynamically adapts the certification horizon between solver iterations.

The architecture is evaluated in Gazebo simulation on assisted feeding and pick-and-place tasks using a Franka Emika Panda and a uFactory xArm7 manipulator. Across 60 unfiltered trials, nominal execution causes 51 collisions, all eliminated by the safety filter. The Panda completes 27 of 30 pick-and-place trials with a median tool deviation of 0.60~cm, while the xArm7 completes 16 of 30 trials without collision. Transfer to the second manipulator requires substituting its kinematic model and retraining only its per-link distance fields. These results show that a decoupled reactive filter can preserve the flexibility of a learned policy while enforcing the learned safety barrier under its stated assumptions.

\end{otherlanguage}
